% ============================================================
%  Příprava na maturitu – Mechatronika
%  Kompiluj: pdflatex maturita_mechatronika.tex
%  Kódování: UTF-8
% ============================================================

\documentclass[a4paper,10pt]{article}

% --- Balíčky ---
\usepackage[T1]{fontenc}
\usepackage[utf8]{inputenc}
\usepackage[czech]{babel}
\usepackage[left=2.4cm, right=1.8cm, top=1.8cm, bottom=1.8cm]{geometry}
\usepackage{lmodern}
\usepackage{microtype}
\usepackage{parskip}
\usepackage{titlesec}
\usepackage{enumitem}
\usepackage{xcolor}
\usepackage{hyperref}
\usepackage{tabularx}
\usepackage{booktabs}
\usepackage{array}
\usepackage{tikz}
\usepackage{eso-pic}
\usepackage{mdframed}
\usepackage{amsmath}
\usepackage{amssymb}
\usepackage{pgfplots}
\pgfplotsset{compat=1.18}
\usetikzlibrary{arrows.meta, positioning, shapes.geometric, calc}

% --- Barvy ---
\definecolor{accent}{HTML}{03254E}
\definecolor{stripeblue}{HTML}{568EA3}
\definecolor{medgray}{HTML}{888888}
\definecolor{lightblue}{HTML}{EAF2F8}
\definecolor{lightyellow}{HTML}{FDFBE4}
\definecolor{lightgreen}{HTML}{EAF7EA}

% --- Modrý pruh vlevo ---
\AddToShipoutPictureBG{%
  \begin{tikzpicture}[remember picture, overlay]
    \fill[stripeblue]
      (current page.north west)
      rectangle
      ([xshift=10mm] current page.south west);
    \fill[accent!60]
      ([xshift=10mm] current page.north west)
      rectangle
      ([xshift=12mm] current page.south west);
  \end{tikzpicture}%
}

\hypersetup{colorlinks=true, urlcolor=accent, linkcolor=accent}

\titleformat{\section}
  {\large\bfseries\color{accent}}
  {\thesection.}
  {0.5em}
  {}
  [\color{accent}\titlerule]

\titlespacing{\section}{0pt}{16pt}{6pt}

\newmdenv[
  backgroundcolor=lightblue,
  linecolor=stripeblue,
  linewidth=1pt,
  roundcorner=4pt,
  innerleftmargin=8pt,
  innerrightmargin=8pt,
  innertopmargin=6pt,
  innerbottommargin=6pt,
  skipabove=4pt,
  skipbelow=4pt
]{pojmybox}

\newmdenv[
  backgroundcolor=lightyellow,
  linecolor=accent!40,
  linewidth=1pt,
  roundcorner=4pt,
  innerleftmargin=8pt,
  innerrightmargin=8pt,
  innertopmargin=6pt,
  innerbottommargin=6pt,
  skipabove=4pt,
  skipbelow=8pt
]{vysvetlenibox}

\newmdenv[
  backgroundcolor=lightgreen,
  linecolor=accent!30,
  linewidth=1pt,
  roundcorner=4pt,
  innerleftmargin=8pt,
  innerrightmargin=8pt,
  innertopmargin=6pt,
  innerbottommargin=6pt,
  skipabove=4pt,
  skipbelow=8pt
]{vzorcebox}

\pagestyle{plain}

% ============================================================
\begin{document}

% --- Záhlaví ---
\begin{minipage}[t]{0.75\textwidth}
    {\Huge\bfseries\color{accent} Příprava na maturitu}\\[4pt]
    {\large\color{stripeblue} Mechatronika}\\[2pt]
    \textcolor{medgray}{\small Straka Jakub, Suchánek Lukáš \quad SPŠ Prosek \quad 2026}
\end{minipage}
\hfill
\begin{minipage}[t]{0.22\textwidth}
    \raggedleft
    \begin{tikzpicture}
        \draw[color=stripeblue, rounded corners=4pt, line width=1pt]
            (0,0) rectangle (3,1.2)
            node[midway, align=center, color=accent, font=\small\bfseries]
            {Otázek: 22};
    \end{tikzpicture}
\end{minipage}

\vspace{6pt}
\noindent\color{accent}\rule{\linewidth}{1.5pt}
\vspace{2pt}

% ============================================================
\section{Převody točivého pohybu (ozubené, řetězové, řemenové)}

\begin{pojmybox}
\textbf{Klíčové pojmy:}
převodový poměr, momentový poměr, účinnost, ozubené kolo,
modul, řemen (plochý, klínový, ozubený), řetěz,
vůle, mazání, evolventní profil, čelní/kuželové ozubení
\end{pojmybox}

\begin{vysvetlenibox}
\textbf{Základní vztahy}

\begin{vzorcebox}
Převodový poměr: $i = \frac{n_1}{n_2} = \frac{z_2}{z_1} = \frac{d_2}{d_1}$\\
Momentový přenos: $M_2 = M_1 \cdot i \cdot \eta$
\quad ($\eta$ = účinnost soukolí)\\
Průměr roztečné kružnice: $d = m \cdot z$
\quad (modul $m$ [mm], počet zubů $z$)\\
Osová vzdálenost čelního soukolí:
$a = \frac{m(z_1 + z_2)}{2}$
\end{vzorcebox}

\textbf{Typy ozubených kol}

\begin{center}
\begin{tabularx}{0.95\linewidth}{@{} l X l @{}}
    \toprule
    \textbf{Typ} & \textbf{Vlastnosti} & \textbf{Účinnost} \\
    \midrule
    Čelní přímé    & Jednoduché, hlučnější, riziko kavitace & 97--99\,\% \\
    Čelní šikmé    & Tišší, větší axiální síla & 97--99\,\% \\
    Kuželové       & Změna osy rotace (90°) & 96--98\,\% \\
    Šneková        & Velký převodový poměr, samosvornost & 50--90\,\% \\
    Planetové      & Kompaktní, velký moment, koaxiální & 95--98\,\% \\
    \bottomrule
\end{tabularx}
\end{center}

\textbf{Řemenové a řetězové převody}

\begin{center}
\begin{tabularx}{0.95\linewidth}{@{} l X X X @{}}
    \toprule
    & \textbf{Klínový řemen} & \textbf{Ozubený řemen}
    & \textbf{Řetěz} \\
    \midrule
    Skluz      & Ano (3--5\,\%) & Ne & Ne \\
    Tlumení    & Vysoké & Střední & Nízké \\
    Mazání     & Ne & Ne & Ano \\
    Účinnost   & 92--96\,\% & 96--99\,\% & 96--98\,\% \\
    Použití    & Kompresory, ventilátory & Rozvod, servopohony
               & Motorky, pohony os \\
    \bottomrule
\end{tabularx}
\end{center}

\textbf{Základní typy převodovek}
\begin{center}
\begin{tabular}{|l|c|c|c|c|}
\hline
 & Běžná (čelní) & Planetární & Harmonická & Cykloidní \\
\hline
Princip & Ozubená kola & Kolo + satelity & Pružná deformace & Excentrický kotouč \\
\hline
Převod & Nízký--střední & Střední--vysoký & Velmi vysoký & Velmi vysoký \\
\hline
Kompaktnost & Nízká & Vysoká & Velmi vysoká & Vysoká \\
\hline
Vůle & Střední & Nízká & Velmi nízká & Velmi nízká \\
\hline
Účinnost & 95--98\,\% & 95--97\,\% & 75--90\,\% & 85--95\,\% \\
\hline
Použití & Stroje & Robotika & Přesné pohony & Reduktory \\
\hline
\end{tabular}
\end{center}

\end{vysvetlenibox}

% ============================================================

% ============================================================
\end{document}
